\section{検証 III：3 箱}\label{sec:three}

3 箱には非長方形の彩色 $\yng(2,1)$ が要る．
定義どおりの $\rCS_{\yng(2,1)}$ が得られるのは，トーラス結び目と iterated torus 結び目（DAHA，楕円 Hall 代数）である．
この節の主な結果は，これらの結び目についてのものである．

\subsection{トーラス結び目と代数的な iterated torus 結び目}
% [R10, R11]
\begin{compresult}\label{cr:three}
表~\ref{tab:three} の 15 個の結び目で，$|\rho|\le3$ の全 6 つの $\rho$ について $\hf_\rho$ は正（framing 0）である．
内訳はトーラス結び目 7 個，代数的な iterated torus 結び目 7 個，$\mathbf{4_1}$ である．
同じ全体の符号で，$G$ 展開の前の $\hbN$ と後の $\hN$ がともに非負整数になる．
\end{compresult}

\begin{table}[h]
\centering
\caption[3 箱までの判定（framing 0）]{3 箱までの判定（framing 0）．$\checkmark(n)$ は $G$ 展開の前後とも負の項が 0 で，$n$ は展開の前の項数．
$T_{2,3}$，$T_{2,5}$ と $\mathbf{4_1}$ は KN の表の再現である．$\mathbf{4_1}$ の $\yng(2,1)$ は KN の候補を用いた．}\label{tab:three}
\resizebox{\textwidth}{!}{% scripts/paper_summary.py が生成。\checkmark(n) は spin 表示と genus 表示の両方で負の項 0（n は spin 表示の項数）。それ以外は spin 表示の負/項数; genus 表示の負/項数（前後それぞれで全体の符号は良い方）。framing 0。
\begin{tabular}{llcccccc}
\hline
結び目 & 分類 & $\yng(1)$ & $\yng(2)$ & $\yng(1,1)$ & $\yng(3)$ & $\yng(2,1)$ & $\yng(1,1,1)$ \\
\hline
$T_{2,3}$ & torus & \checkmark\,(6) & \checkmark\,(26) & \checkmark\,(41) & \checkmark\,(128) & \checkmark\,(236) & \checkmark\,(282) \\
$T_{2,5}$ & torus & \checkmark\,(10) & \checkmark\,(112) & \checkmark\,(144) & \checkmark\,(753) & \checkmark\,(986) & \checkmark\,(1083) \\
$T_{3,4}$ & torus & \checkmark\,(21) & \checkmark\,(429) & \checkmark\,(504) & \checkmark\,(3249) & \checkmark\,(3802) & \checkmark\,(4053) \\
$T_{3,5}$ & torus & \checkmark\,(30) & \checkmark\,(766) & \checkmark\,(867) & \checkmark\,(5831) & \checkmark\,(6559) & \checkmark\,(6890) \\
$T_{2,11}$ & torus & \checkmark\,(22) & \checkmark\,(733) & \checkmark\,(813) & \checkmark\,(5667) & \checkmark\,(6260) & \checkmark\,(6519) \\
$T_{3,7}$ & torus & \checkmark\,(53) & \checkmark\,(1797) & \checkmark\,(1953) & \checkmark\,(13600) & \checkmark\,(14696) & \checkmark\,(15199) \\
$T_{4,5}$ & torus & \checkmark\,(68) & \checkmark\,(2356) & \checkmark\,(2558) & \checkmark\,(18729) & \checkmark\,(20195) & \checkmark\,(20883) \\
K1 & alg & \checkmark\,(104) & \checkmark\,(4141) & \checkmark\,(4414) & \checkmark\,(32732) & \checkmark\,(34655) & \checkmark\,(35558) \\
K3 & alg & \checkmark\,(122) & \checkmark\,(5178) & \checkmark\,(5483) & \checkmark\,(41138) & \checkmark\,(43289) & \checkmark\,(44300) \\
K6a & alg & \checkmark\,(275) & \checkmark\,(12784) & \checkmark\,(13267) & \checkmark\,(100115) & \checkmark\,(103458) & \checkmark\,(105035) \\
K9 & alg & \checkmark\,(564) & \checkmark\,(28584) & \checkmark\,(29309) & \checkmark\,(224004) & \checkmark\,(228995) & \checkmark\,(231354) \\
K8 & alg & \checkmark\,(617) & \checkmark\,(31499) & \checkmark\,(32427) & \checkmark\,(251627) & \checkmark\,(258022) & \checkmark\,(261098) \\
K7 & alg & \checkmark\,(834) & \checkmark\,(43533) & \checkmark\,(44631) & \checkmark\,(343739) & \checkmark\,(351218) & \checkmark\,(354818) \\
K10 & alg & \checkmark\,(4119) & \checkmark\,(227727) & \checkmark\,(230642) & \checkmark\,(1805429) & \checkmark\,(1825082) & \checkmark\,(1834640) \\
$4_1$ & kn & \checkmark\,(10) & \checkmark\,(42) & \checkmark\,(42) & \checkmark\,(223) & \checkmark\,(264) & \checkmark\,(223) \\
\hline
\end{tabular}
}
\end{table}

$T_{2,3}$，$T_{2,5}$ と $\mathbf{4_1}$ の結果は，KN の表の再現である．
KN の表に無いトーラス結び目 5 個（$T_{3,4}$ の 3 箱，$T_{3,5}$，$T_{2,11}$，$T_{3,7}$，$T_{4,5}$）は新しい表である．
ただし KN はトーラス結び目全般について予想しているので，これは予想の確認に当たる．
代数的な iterated torus 結び目 7 個の 3 箱は，本論文の調べた範囲では文献に無い．
3 本ケーブル $K7$ と 2 段の反復 $K10$ では，$|\lambda|=3$ の refined 不変量そのものも文献に見当たらない．
Cherednik--Danilenko~\cite{Cherednik14} は，$K7$ の $\yng(1)$ と $K10$ の DAHA-Betti 多項式を与えている．

% [R12]
\begin{remark}
入力の $\rCS_{\yng(2,1)}$ 自体は負の係数をもつが，$\hf_{\yng(2,1)}$ は正になる．
したがって，3 箱の正値性は入力の係数の正値性から直接は従わず，KN の抽出（plethystic な対数と $M$ 行列）を経て初めて現れる．
\end{remark}

\subsection{framing を変えた 3 箱}
% [R13]
\begin{compresult}\label{cr:three-x}
表~\ref{tab:three} の全ての結び目について，3 箱を $x\neq0$ でも判定した．
トーラス結び目 $T_{3,5},T_{2,11},T_{3,7},T_{4,5}$ と代数的なケーブル $K6a,K7,K8,K9$ は $|x|\le10$，$K10$ は $|x|\le2$，$K1$ と $K3$ は $|x|\le2$（$K1$ は $x=\pm6$ も），
自明な結び目と $T_{2,3}$ は $|x|\le6$，$T_{2,5}$ と $\mathbf{4_1}$ は $|x|\le4$，$T_{2,7}$ と $T_{3,4}$ は $|x|\le3$ の範囲である．
いずれも正である（計算結果~\ref{cr:framing-batch} の一部）．
$x$ を正に大きくすると項数が増え，負では減る（$K10$ の $\rho=\yng(3)$ で，$x=-2$ では約 175 万項，$x=2$ では約 186 万項）．
\end{compresult}
3 箱でも $\bt^{-3x}\hf^{(x)}_\rho$ は 7 根の指数和である（命題~\ref{prop:expsum}）．
しかし，6 根のニュートン形では一様な符号の並びが見つからず，全ての $x$ の証明書は得られていない．

\subsection{予想 B とその動機}\label{sec:conjB}
% [R14]
\begin{tcolorbox}
\begin{conjecture}\label{conj:B}
代数的な iterated torus 結び目は，framing 0 で全ての $\rho$ について正である．
\end{conjecture}
\end{tcolorbox}

KN の物理的な議論は，追加の $U(1)_R$ 対称性を Seifert 多様体の構造から得る．
代数的な iterated torus 結び目の補空間は一般に Seifert 多様体ではないので，この議論は直接は当てはまらない．
予想~\ref{conj:B} の動機は，次の Cherednik--Danilenko~\cite{Cherednik14} の予想にある．
\begin{itemize}
\item 正値性の予想：代数的な結び目と任意の長方形のヤング図形について，DAHA superpolynomial の係数は全て正である（このような結び目を pseudo-algebraic と呼ぶ）．
\item Connection Conjecture (ii)：彩色なしの pseudo-algebraic な結び目では，DAHA superpolynomial は（tilde 正規化した）reduced stable Khovanov--Rozansky 多項式に一致する．
\end{itemize}
これらと本論文の結果の関係は次のとおりである．
\begin{itemize}
\item $\yng(1)$：Connection Conjecture (ii) の下で，DAHA の $\yng(1)$ は KhR の Poincaré 多項式 $P_{\mathrm{KhR}}$ の単項式倍 $mP_{\mathrm{KhR}}$ になる．
(N1) により両者は $\bt=-1$ で同じ reduced HOMFLY-PT 多項式を与えるので，$m=\pm\bt^k$ で $\pm(-1)^k=1$ である．両者とも係数が非負なので $m=\bt^{2l}$ である．
$\bt^{2l}$ は $\sigma$ で不変なので，命題~\ref{prop:one} が適用でき，$\yng(1)$ の正値性が従う．
\item 長方形：Cherednik--Danilenko は DAHA の値の正値性を予想し，ホモロジーとの一致を期待している．予想~\ref{conj:B} の長方形の部分は，これを KN の BPS の段に移したものである．
\item $\yng(2,1)$：Cherednik--Danilenko は長方形でない彩色の正値性を主張していない．実際 $\rCS_{\yng(2,1)}$ は負の係数をもつ．BPS の段での正値性は，その範囲の外の現象である．
\item 非代数的なケーブル：Cherednik--Danilenko は，正値性をもたない非代数的な iterated torus 結び目があることを認めている．本論文の非代数的なケーブルの多くで，DAHA の値の $\yng(1)$ から負の係数が現れる（付録~\ref{app:nonalg}）．
\end{itemize}

\subsection{双曲結び目の 3 箱へ：非長方形の彩色の問題}\label{sec:21}
双曲結び目では，定義に基づく refined $\yng(2,1)$ が文献にもソフトウェアにも無い．
そこで，構造的な条件から $\yng(2,1)$ を決められるかを調べた．
% [R15, R16, R17]
\begin{itemize}
\item 結び目ごと：条件は $t=q$ での 3 本ケーブルから得た $\yng(2,1)$ 彩色 HOMFLY-PT 多項式，ミラー，指数増大性，$U(2)$ への簡約，3 箱の Laurent 性，$G$ 展開である．
これらを課しても $\yng(2,1)$ は一意に決まらない．解の差は「結び目によらない因子 × 自由な多項式」の形をしている．
曖昧さがあること自体は Anokhina ら~\cite{Anokhina12} も指摘している．
\item 弱い判定：正値性を満たす整数の $\yng(2,1)$ は，$\mathbf{4_1}$（KN の候補）と $\mathbf{6_1}$（813 項の解）に存在する．$\mathbf{5_2}$ では有理数の解しか見つからない（整数解の有無は未決）．
これは予想の検証ではなく，予想と矛盾しないことの確認にとどまる．
\item 族として：$t=q$ ではダブルツイスト結び目の $\yng(2,1)$ について Morozov の式~\cite{Morozov16,Morozov19} がケーブルからの値と一致する．これは既知の式の検証である．
refined の長方形は KNTZ~\cite{Kameyama19} の式を実装し，GNSSS の式と三葉結び目の楕円 Hall 代数の値（$\yng(2,2)$ まで）と照合した．しかし refined の $\yng(2,1)$ は，族の構造の仮定を置いても一意に決まらない．
\end{itemize}
双曲結び目の 3 箱には，定義に基づく $\yng(2,1)$ が必要である．その候補は \S\ref{sec:discussion} で述べる．
